\section{Introduction}
Large language model (LLM) agents are increasingly expected to operate over long horizons: they must remember user preferences, track environment changes, integrate tool outputs, and use this information to answer future questions, plan, and act.
A coding assistant told that a project has moved from REST to GraphQL should not preserve both as equally live assumptions; a data-analysis agent told to ignore a corrupted dataset should revise downstream conclusions; an assistant asked to forget a user preference should update future behavior immediately. These are not primarily failures of recall. They are failures of \emph{state maintenance}.

Large context windows and retrieval-augmented generation (RAG) address part of the memory problem by making more history accessible at inference time. But access to history is not the same as maintaining a usable current state. Recent benchmarks confirm this gap. LoCoMo and LongMemEval show that long-range temporal reasoning and knowledge revision remain difficult even for long-context and retrieval-augmented systems \cite{maharana2024locomo,wu2025longmemeval}. MemoryArena pushes the problem further: systems that perform well on recall-oriented benchmarks remain brittle when memory must support future action in interdependent multi-session agentic tasks \cite{he2026memoryarena}. The bottleneck is not whether past information can be retrieved, but whether interaction history can be consolidated into a compact state that remains useful for future decisions.

We argue that most existing approaches remain \emph{archive-first}, not state-first. Long-context methods preserve raw history; RAG and note-style memory rewrite history into retrievable chunks, summaries, or facts; more explicit memory systems such as Mem0, MemSkill, and UMA improve memory writing and maintenance, but still largely center the storage, update, and retrieval of memory items \cite{chhikara2025mem0,zhang2026memskill,zhang2026uma}. As a result, current state is often reconstructed at decision time by rereading context or retrieved notes, rather than maintained as a first-class object throughout interaction. Recent analysis further suggests that retrieval choices can matter more than sophisticated write-time memory construction, reinforcing the view that many current systems learn a better archive, not a better state \cite{yuan2026diagnosing}.

\paragraph{Motivating example.}
Consider a project-management stream for ticket \texttt{bug\#512}: Alice opens it on Day~1 with blocker \emph{schema migration}; Bob takes over on Day~3 and adds \emph{auth error} as a second blocker; Alice reports the schema migration resolved on Day~5; Carol takes over on Day~7; on Day~9 Carol closes the auth error but reports a new blocker, \emph{rate-limit timeout}. A downstream query asks: \emph{Who owns \texttt{bug\#512}, and what are the open blockers?}
The correct answer (Carol; rate-limit timeout) is not a retrievable span. A retrieval system will surface Days~3, 7, and~9 based on keyword overlap with ``blocker,'' but the agent must still reconcile at query time which updates supersede which. Value-aware retrieval~\citep{zhang2026memrl} improves \emph{which} entries are selected but does not resolve contradictions among them; CRUD-style memory banks~\citep{zhang2026uma} address this for summarization but lack a learned selection mechanism. The task requires maintaining an editable \emph{current state}, not selecting from a historical archive.

In this paper, we formulate agent memory as \emph{budgeted adaptive state tracking}. Given a growing interaction history $x_{1:t}$ and a memory budget $B$, the agent maintains a memory substrate $z_t = f_{\theta}(x_{1:t}; B)$ that is compact enough to fit the budget while remaining maximally useful for future answering, planning, tool use, and correction. Each entry in $z_t$ carries both a utility estimate (how valuable it is for future decisions) and a belief state (how reliable it is given observed outcomes). Each new interaction induces an explicit state transition on $z_t$: \textsc{Create}, \textsc{Refine}, \textsc{Override}, \textsc{Merge}, \textsc{Expire}, \textsc{Delete}, or \textsc{Redact}, rather than append-only note accumulation. This makes contradiction resolution, temporal overwriting, local editing, and controllable deletion first-class operations. Downstream policies operate \emph{state-first}: given a task query, the system compiles an executable state view $\hat{s}_t(q)$ from $z_t$---selecting relevant entries, filtering expired material, and organizing results by confidence---so that the agent reads a pre-processed operating state rather than raw history or unscored candidates. Two properties follow: the state view is \emph{executable}, directly supporting downstream decisions rather than serving as a retrieval index, and \emph{intervention-ready}, so that user corrections, policy changes, or new evidence immediately affect subsequent behavior through the operator interface.

\paragraph{Contributions.}
Our main contributions are threefold. First, we recast agent memory as budgeted adaptive state tracking and identify maintained operating state, rather than archival storage, as the central object for long-horizon agents. Second, we introduce explicit state-transition operators with dual-rate feedback (Q-learning for utility, Beta posteriors for belief), enabling consistency maintenance, conflict-driven rewriting, and budget-aware retention on an episodic memory substrate. Third, we propose a state-first execution interface---where the agent reads a compiled, confidence-organized state view rather than raw history---and evaluation axes centered on state maintenance (consistency, intervention fidelity, conflict resolution, budget-utility trade-offs) that are not captured by recall-oriented memory benchmarks alone.
